\section{Background \label{sec:background}}

Software artifacts are naturally relational and multi-typed. A system contains folders, source files or classes, and finer-grained program elements such as methods and fields, connected through relations that include folder containment, file-to-member composition, and static code-level dependencies.

To preserve these distinctions, we represent each software system as a typed graph:
$$G=(V, E,\mathcal{A},\mathcal{R}),$$
where $V$ and $E$ are the sets of artifact nodes and relations between them, respectively. The mappings $\tau: V \rightarrow \mathcal{A}$ and $\phi: E \rightarrow \mathcal{R}$ assign node and relation types. In our most detailed configuration, $\mathcal{A}$ contains folder, file, and aggregate member types, while $\mathcal{R}$ distinguishes folder containment, file-to-member composition, and a single dependency relation that merges all file-level static dependencies. The file-level graph is homogeneous, and adding folder or member nodes makes it heterogeneous.

A relation's meaning depends on its type and the types of its source and target nodes in $G$. We describe an edge $e_{s,t}$ from node $s$ to node $t$ by the meta-relation $\langle \tau(s), \phi(e_{s,t}), \tau(t) \rangle$. For example, $\langle \text{folder}, \text{contains}, \text{file} \rangle$ captures project organization, whereas $\langle \text{file}, \text{depends on}, \text{file} \rangle$ captures code-level coupling. These meta-relations preserve the distinction between structural organization and implementation dependencies, which a homogeneous graph with one node type and one relation type does not explicitly represent.

HGT~\cite{hgt} exploits such meta-relations during message passing by conditioning attention and message transformation on the source node type, edge type, and target node type. GAT~\cite{gat} also uses attention-based message passing, but its standard formulation does not explicitly distinguish node or relation types. HGT can also operate on a graph with a single node type and a single relation type, although this does not make its attention mechanism equivalent to that of GAT.

The initial node features determine what information is available before message passing, while graph structure determines which nodes exchange information. Depending on how they are encoded, a file's features can capture lexical evidence from its identifiers, structural evidence from its position in the dependency graph, or contextual evidence from its code.

A lexical embedding trained on identifier tokens captures naming patterns but does not explicitly encode dependency structure or implementation syntax and behavior. A graph embedding such as Node2Vec~\cite{node2vec} instead encodes structural neighborhoods without using code content. Such embeddings have been used to capture dependency-based cohesion in architecture recovery~\cite{negar,structsar}.

Contextual representations from pretrained code encoders such as UniXcoder~\cite{unixcoder} capture the syntactic and behavioral content of an implementation that lexical features miss, but their pretraining objectives do not target architectural grouping: files with similar code can belong to different modules. Contrastive learning adapts an encoder to a chosen notion of similarity by bringing designated positive pairs closer than other in-batch examples treated as negatives, using an InfoNCE objective~\cite{infonce}. Dependency-derived positive pairs encourage embeddings to reflect dependency proximity, which need not align with reference-module membership.